\documentclass[11pt,a4paper]{article}

% ---------- Encodage et langue ----------
\usepackage[utf8]{inputenc}
\usepackage[T1]{fontenc}
\usepackage{lmodern}
\usepackage[french]{babel}
\usepackage{microtype}

% ---------- Mise en page ----------
\usepackage[a4paper, margin=2.5cm, headheight=24pt, footskip=1.2cm]{geometry}
\usepackage{xcolor}
\usepackage{graphicx}
\usepackage{array}
\usepackage{float}
\usepackage{tikz}
\usetikzlibrary{calc}
\usepackage{fancyhdr}
\usepackage{titlesec}
\usepackage{caption}
\usepackage{listings}
\usepackage[hidelinks]{hyperref}

\linespread{1.1}
\setlength{\parindent}{0pt}
\setlength{\parskip}{6pt}

% ---------- Couleurs ENIAD ----------
\definecolor{ENIADNavy}{HTML}{12304A}
\definecolor{ENIADOrange}{HTML}{E87924}
\definecolor{ENIADBlue}{HTML}{2878A8}
\definecolor{ENIADGray}{HTML}{5D6872}
\definecolor{ENIADWhite}{HTML}{FFFFFF}
\definecolor{ENIADLightBlue}{HTML}{EAF3F8}
\definecolor{CodeGreen}{HTML}{3C7A4E}

% ---------- En-tête et pied de page ----------
\pagestyle{fancy}
\fancyhf{}
\renewcommand{\headrulewidth}{0.8pt}
\renewcommand{\footrulewidth}{0pt}
\renewcommand{\headrule}{\hbox to\headwidth{\color{ENIADOrange}\leaders\hrule height \headrulewidth\hfill}}
\fancyhead[L]{\small\sffamily\color{ENIADGray} Big Data \textbullet\ TP\,2}
\fancyhead[R]{\small\sffamily\color{ENIADGray} Zakariae El Haddouchi}
\fancyfoot[C]{\small\sffamily\color{ENIADNavy}\thepage}

% ---------- Titres de section ----------
\titleformat{\section}[block]
  {\sffamily\Large\bfseries\color{ENIADNavy}}
  {\textcolor{ENIADOrange}{Exercice~\thesection}\hspace{0.5em}\textcolor{ENIADOrange}{|}}
  {0.6em}
  {}
  [\vspace{3pt}{\color{ENIADOrange}\titlerule[1pt]}]
\titlespacing*{\section}{0pt}{20pt}{10pt}

% ---------- Légendes ----------
\addto\captionsfrench{\renewcommand{\figurename}{Fig.}}
\captionsetup{
  font={small,rm,color=ENIADGray},
  labelfont={up,color=ENIADOrange},
  labelsep=colon,
  skip=6pt
}

% ---------- Code ----------
\lstdefinestyle{eniad}{
  language=Python,
  basicstyle=\ttfamily\footnotesize,
  keywordstyle=\color{ENIADBlue}\bfseries,
  stringstyle=\color{ENIADOrange},
  commentstyle=\color{CodeGreen}\itshape,
  numbers=left,
  numberstyle=\tiny\color{ENIADGray},
  numbersep=8pt,
  backgroundcolor=\color{ENIADLightBlue!45},
  frame=leftline,
  rulecolor=\color{ENIADOrange},
  framerule=2pt,
  framesep=6pt,
  xleftmargin=22pt,
  framexleftmargin=18pt,
  breaklines=true,
  showstringspaces=false,
  columns=fullflexible,
  aboveskip=4pt,
  belowskip=8pt,
  literate=
    {é}{{\'e}}1 {è}{{\`e}}1 {ê}{{\^e}}1 {à}{{\`a}}1 {â}{{\^a}}1
    {ç}{{\c c}}1 {î}{{\^\i}}1 {ô}{{\^o}}1 {û}{{\^u}}1 {ù}{{\`u}}1
    {É}{{\'E}}1 {’}{{'}}1
}
\lstset{style=eniad}

% Nom de fichier affiché au-dessus d'un bloc de code
\newcommand{\fichier}[1]{%
  \par\medskip\noindent
  {\sffamily\small\bfseries\color{ENIADNavy}%
   \colorbox{ENIADLightBlue}{\strut\ \texttt{#1}\ }}%
  \par\nopagebreak\vspace{-2pt}}

% Capture d'écran (cadre de remplacement si l'image est absente)
\newcommand{\rapportcapture}[2]{%
  \begin{figure}[H]
    \centering
    \IfFileExists{#1}{%
      \fbox{\includegraphics[width=0.9\linewidth]{#1}}%
    }{%
      \fcolorbox{ENIADBlue!50}{ENIADLightBlue}{%
        \parbox[c][4cm][c]{0.82\linewidth}{\centering\sffamily\color{ENIADGray}
          Capture à insérer\\[2pt]{\footnotesize\ttfamily\detokenize{#1}}}}%
    }
    \caption{#2}
  \end{figure}}

\title{TP2}
\author{}
\date{}

\begin{document}

% Required in preamble:
% \usepackage{tikz}  \usetikzlibrary{calc}
% \usepackage{xcolor}  \usepackage{graphicx}  \usepackage{array}
% \usepackage[T1]{fontenc}  \usepackage{lmodern}   % needed for the large title size
%
% Colors (skip if already defined; \providecolor won't overwrite):
\providecolor{ENIADNavy}{HTML}{0B1F3A}
\providecolor{ENIADBlue}{HTML}{1F4E8C}
\providecolor{ENIADLightBlue}{HTML}{EEF3FA}
\providecolor{ENIADOrange}{HTML}{F28C28}
\providecolor{ENIADGray}{HTML}{5B6573}
\providecolor{ENIADWhite}{HTML}{FFFFFF}

\begin{titlepage}
\thispagestyle{empty}
\sffamily

\begin{tikzpicture}[remember picture, overlay]

  % ---------- Background ----------
  \fill[ENIADWhite] (current page.south west) rectangle (current page.north east);

  % subtle decorative circles (top-right of body, bottom-left)
  \fill[ENIADLightBlue] ([xshift=1cm,yshift=-6cm] current page.north east) circle (4.5cm);
  \fill[ENIADLightBlue!60] ([xshift=-1cm,yshift=5cm] current page.south west) circle (3.5cm);

  % ---------- Header band ----------
  \fill[ENIADNavy]
    (current page.north west) rectangle ([yshift=-4.2cm] current page.north east);
  \fill[ENIADOrange]
    ([yshift=-4.2cm] current page.north west) rectangle ([yshift=-4.45cm] current page.north east);

  \node[anchor=west, text=ENIADWhite, align=left, text width=11cm]
    at ([xshift=2cm, yshift=-2.1cm] current page.north west)
    {{\Large\bfseries Université Mohammed Premier}\\[5pt]
     {\small École Nationale de l'Intelligence Artificielle\\et du Digital --- Berkane}};

  \node[anchor=east, inner sep=0pt]
    at ([xshift=-2cm, yshift=-2.1cm] current page.north east)
    {\IfFileExists{logo-eniad.png}
      {\includegraphics[height=2.8cm]{logo-eniad.png}}
      {\textcolor{ENIADWhite}{\Huge\bfseries ENIAD}}};

  % ---------- Title block ----------
  \node[anchor=north, text=ENIADOrange, font=\bfseries\small, align=center]
    at ([yshift=-6.6cm] current page.north)
    {C O M P T E \ \ R E N D U \ \ D E \ \ T R A V A U X \ \ P R A T I Q U E S};

  \node[anchor=north, text=ENIADNavy, font=\fontsize{34}{40}\selectfont\bfseries, align=center]
    at ([yshift=-7.7cm] current page.north)
    {BIG DATA};

  \fill[ENIADOrange, rounded corners=1pt]
    ([xshift=-1.5cm, yshift=-9.9cm] current page.north)
    rectangle ([xshift=1.5cm, yshift=-10.05cm] current page.north);

  \node[anchor=north, text=ENIADBlue, font=\Large\bfseries, align=center, text width=15cm]
    at ([yshift=-10.6cm] current page.north)
    {Cluster HDFS avec Docker Compose};

  % ---------- TP badge ----------
  \node[anchor=north, fill=ENIADNavy, text=ENIADWhite, font=\bfseries\Large,
        minimum width=3.6cm, minimum height=1.1cm, rounded corners=5pt]
    at ([yshift=-13.2cm] current page.north)
    {TP\,2};

  % ---------- Info card ----------
  \node[anchor=north, name=card, fill=ENIADLightBlue, inner xsep=22pt, inner ysep=16pt,
        text width=11.8cm, align=left]
    at ([yshift=-15.6cm] current page.north)
    {%
      \renewcommand{\arraystretch}{1.7}%
      \begin{tabular}{@{}p{3.3cm}p{8.0cm}@{}}
        {\color{ENIADGray}\footnotesize\bfseries RÉALISÉ PAR} &
          {\color{ENIADNavy}\large\bfseries Zakariae El Haddouchi} \\
        {\color{ENIADGray}\footnotesize\bfseries FILIÈRE} &
          {\color{ENIADNavy}\normalsize Intelligence Artificielle (cycle ingénieur)} \\
        {\color{ENIADGray}\footnotesize\bfseries MODULE} &
          {\color{ENIADNavy}\normalsize Big Data} \\
        {\color{ENIADGray}\footnotesize\bfseries PROFESSEUR} &
          {\color{ENIADNavy}\normalsize\bfseries Phd. AL ALAMI NABIL} \\
      \end{tabular}%
    };
  \fill[ENIADOrange] (card.north west) rectangle ([xshift=0.18cm] card.south west);

  % ---------- Footer ----------
  \fill[ENIADNavy]
    (current page.south west) rectangle ([yshift=1.8cm] current page.south east);
  \fill[ENIADOrange]
    ([yshift=1.8cm] current page.south west) rectangle ([yshift=1.95cm] current page.south east);
  \node[text=ENIADWhite, font=\small\bfseries, anchor=center]
    at ([yshift=0.9cm] current page.south)
    {Année universitaire 2026\,--\,2027};

  \node[anchor=north, align=center, text=ENIADNavy]
    at ([yshift=-23.5cm] current page.north)
    {{\footnotesize\bfseries CODE SOURCE}\\[3pt]
     {\footnotesize\href{https://github.com/Eng-ZAKARIAE-AI/bigdata-hadoop}
       {github.com/Eng-ZAKARIAE-AI/bigdata-hadoop}}};

\end{tikzpicture}

\end{titlepage}
\clearpage

\section{Présentation du cluster HDFS}

Ce TP vise à déployer un cluster Hadoop HDFS minimal (1 NameNode + 2 DataNodes) en utilisant Docker Compose. Le cluster repose sur l'image officielle \texttt{apache/hadoop:3.5.0}, avec une configuration personnalisée montée en volume.

\subsection{Structure du projet}

\begin{itemize}
  \item \texttt{Dockerfile}
  \item \texttt{docker-compose.yml}
  \item \texttt{config/} : \texttt{core-site.xml}, \texttt{hdfs-site.xml}
  \item \texttt{data/} : \texttt{namenode\_name/}, \texttt{datanode1/}, \texttt{datanode2/}
  \item \texttt{run-hdfs.sh}
  \item \texttt{starter.sh}
  \item \texttt{hadoop-wrapper.sh}
  \item \texttt{README.md}
\end{itemize}

\paragraph{Description des composants.}
\begin{itemize}
  \item \texttt{Dockerfile} : construit l'image \texttt{tp2-hadoop:3.5.0} à partir de \texttt{apache/hadoop:3.5.0}, copie les fichiers de configuration Hadoop et définit \texttt{hdfs} comme point d'entrée.
  \item \texttt{docker-compose.yml} : orchestre les trois services \texttt{namenode}, \texttt{datanode1} et \texttt{datanode2} sur le réseau \texttt{hdfs-net}.
  \item \texttt{config/core-site.xml} : définit le système de fichiers par défaut (\texttt{fs.defaultFS} = \texttt{hdfs://namenode:8020}).
  \item \texttt{config/hdfs-site.xml} : configure la réplication, la taille des blocs et les répertoires de stockage.
  \item \texttt{data/} : répertoires locaux montés dans les conteneurs pour persister les métadonnées et les blocs HDFS.
\end{itemize}


\section{Configuration HDFS}

\fichier{config/core-site.xml}
\begin{lstlisting}[language=XML]
<?xml version="1.0" encoding="UTF-8"?>
<configuration>
  <property>
    <name>fs.defaultFS</name>
    <value>hdfs://namenode:8020</value>
  </property>
</configuration>
\end{lstlisting}

\fichier{config/hdfs-site.xml}
\begin{lstlisting}[language=XML]
<?xml version="1.0" encoding="UTF-8"?>
<configuration>
  <property>
    <name>dfs.replication</name>
    <value>2</value>
  </property>
  <property>
    <name>dfs.blocksize</name>
    <value>8388608</value>
  </property>
  <property>
    <name>dfs.namenode.name.dir</name>
    <value>/data/hdfs/name</value>
  </property>
  <property>
    <name>dfs.datanode.data.dir</name>
    <value>/data/hdfs/data</value>
  </property>
</configuration>
\end{lstlisting}

\paragraph{Interprétation des paramètres.}
\begin{itemize}
  \item \texttt{dfs.replication = 2} : chaque bloc est répliqué sur 2 DataNodes, ce qui correspond au nombre de DataNodes du cluster.
  \item \texttt{dfs.blocksize = 8\,388\,608} : la taille d'un bloc HDFS est fixée à 8\,Mio (au lieu des 128\,Mio par défaut), ce qui facilite les tests pédagogiques avec de petits fichiers.
  \item \texttt{dfs.namenode.name.dir = /data/hdfs/name} : les métadonnées du NameNode (namespace, images fsimage, edit logs) sont persistées dans ce répertoire monté depuis l'hôte.
  \item \texttt{dfs.datanode.data.dir = /data/hdfs/data} : les blocs de données sont stockés dans ce répertoire sur chaque DataNode.
  \item \texttt{fs.defaultFS = hdfs://namenode:8020} : l'URI HDFS par défaut pointe vers le NameNode sur le port RPC 8020.
\end{itemize}


\section{Docker Compose : orchestration des services}

\fichier{docker-compose.yml}
\begin{lstlisting}
services:
  namenode:
    image: tp2-hadoop:3.5.0
    container_name: tp2-namenode
    hostname: namenode
    user: hadoop
    command: ["namenode"]
    ports:
      - "9870:9870"   # UI web NameNode
      - "8020:8020"   # Port RPC
    volumes:
      - ./config:/opt/hadoop/etc/hadoop:ro
      - ./data/namenode_name:/data/hdfs/name
    networks: [hdfs-net]

  datanode1:
    image: tp2-hadoop:3.5.0
    container_name: tp2-datanode1
    hostname: datanode1
    user: hadoop
    command: ["datanode"]
    depends_on: [namenode]
    volumes:
      - ./config:/opt/hadoop/etc/hadoop:ro
      - ./data/datanode1:/data/hdfs/data
    networks: [hdfs-net]

  datanode2:
    image: tp2-hadoop:3.5.0
    container_name: tp2-datanode2
    hostname: datanode2
    user: hadoop
    command: ["datanode"]
    depends_on: [namenode]
    volumes:
      - ./config:/opt/hadoop/etc/hadoop:ro
      - ./data/datanode2:/data/hdfs/data
    networks: [hdfs-net]

networks:
  hdfs-net:
\end{lstlisting}

\paragraph{Analyse de l'orchestration.}
Le service \texttt{namenode} expose deux ports : \texttt{9870} pour l'interface web d'administration HDFS et \texttt{8020} pour les appels RPC des clients et DataNodes. Les deux DataNodes dépendent du NameNode (\texttt{depends\_on}) et partagent le même réseau Docker \texttt{hdfs-net}, ce qui permet la résolution DNS mutuelle entre les conteneurs (\texttt{namenode}, \texttt{datanode1}, \texttt{datanode2}).

Le montage \texttt{./config:/opt/hadoop/etc/hadoop:ro} écrase les fichiers de configuration inclus dans l'image Docker par ceux du projet, garantissant que le cluster utilise bien notre configuration personnalisée. Les répertoires \texttt{./data/*} assurent la persistance des métadonnées et des blocs entre les redémarrages.


\section{Dockerfile : construction de l'image}

\fichier{Dockerfile}
\begin{lstlisting}
FROM apache/hadoop:3.5.0

COPY config/core-site.xml /opt/hadoop/etc/hadoop/core-site.xml
COPY config/hdfs-site.xml /opt/hadoop/etc/hadoop/hdfs-site.xml

ENTRYPOINT ["/opt/hadoop/bin/hdfs"]
\end{lstlisting}

\paragraph{Explication.}
L'image de base \texttt{apache/hadoop:3.5.0} fournit une installation complète de Hadoop. Les deux fichiers de configuration sont copiés dans le répertoire \texttt{/opt/hadoop/etc/hadoop/} de l'image, remplaçant les fichiers par défaut. L'écrasement du point d'entrée par \texttt{/opt/hadoop/bin/hdfs} permet de contourner le script \texttt{envtoconf.py} défaillant de l'image officielle : chaque conteneur reçoit directement l'argument \texttt{namenode} ou \texttt{datanode} via la commande \texttt{hdfs}.


\section{Démarrage et vérification du cluster}

\subsection{Construction de l'image}

\begin{lstlisting}[language=bash]
docker build -t tp2-hadoop:3.5.0 .
\end{lstlisting}

\subsection{Formatage du NameNode (initialisation uniquement)}

\begin{lstlisting}[language=bash]
docker compose run --rm namenode namenode -format
\end{lstlisting}

Cette étape crée les métadonnées initiales du NameNode dans \texttt{data/namenode\_name/}. Elle n'est à exécuter qu'une seule fois, lors de la première initialisation du cluster.

\subsection{Démarrage du cluster}

\begin{lstlisting}[language=bash]
docker compose up -d
\end{lstlisting}

\rapportcapture
  {screenshots/cluster-up.png}
  {Démarrage du cluster HDFS avec Docker Compose}

\subsection{Vérification de l'état}

\begin{lstlisting}[language=bash]
docker compose ps
docker compose logs -f namenode
docker compose logs -f datanode1 datanode2
\end{lstlisting}

L'interface web du NameNode est accessible à l'adresse \texttt{http://localhost:9870}.


\section{Opérations HDFS courantes}

\subsection{Lister le répertoire racine HDFS}

\begin{lstlisting}[language=bash]
docker compose exec namenode hdfs dfs -ls /
\end{lstlisting}

\rapportcapture
  {screenshots/hdfs-ls-root.png}
  {Liste du répertoire racine HDFS}

\subsection{Créer un répertoire et y déposer un fichier}

\begin{lstlisting}[language=bash]
docker compose exec namenode hdfs dfs -mkdir -p /tp2/data
echo "Bonjour HDFS depuis Docker Compose" > local.txt
docker compose exec namenode hdfs dfs -put local.txt /tp2/data/
\end{lstlisting}

\subsection{Vérifier la réplication des blocs}

\begin{lstlisting}[language=bash]
docker compose exec namenode hdfs fsck /tp2/data/local.txt -files -blocks
\end{lstlisting}

\rapportcapture
  {screenshots/hdfs-fsck.png}
  {Vérification de la réplication des blocs avec fsck}

\paragraph{Interprétation.}
La commande \texttt{fsck} confirme que le fichier \texttt{local.txt} a bien été découpé en blocs et répliqué sur les deux DataNodes conformément à \texttt{dfs.replication = 2}. Chaque bloc apparaît avec deux répliques, ce qui valide le bon fonctionnement de la configuration HDFS du cluster.


\section{Arrêt du cluster}

\begin{lstlisting}[language=bash]
docker compose down
\end{lstlisting}

Cette commande arrête et supprime les conteneurs, mais les données persistent dans les répertoires locaux \texttt{data/} montés en volume. Un redémarrage avec \texttt{docker compose up -d} restaurera le cluster sans nécessiter de nouveau formatage.


\section{Conclusion}

Ce TP a permis de mettre en pratique les concepts fondamentaux de Hadoop HDFS à travers un déploiement Docker Compose :
\begin{itemize}
  \item Comprendre l'architecture NameNode / DataNode et le rôle de chaque composant.
  \item Configurer HDFS via \texttt{core-site.xml} et \texttt{hdfs-site.xml} (réplication, taille de bloc, répertoires de stockage).
  \item Orchestrer un cluster multi-nœuds avec Docker Compose (réseau partagé, volumes de persistance, dépendances entre services).
  \item Exécuter des opérations HDFS courantes (\texttt{mkdir}, \texttt{put}, \texttt{ls}, \texttt{fsck}) depuis le conteneur NameNode.
\end{itemize}

\end{document}